\section{结论}
本章针对状态观测或动作执行中的常值延迟问题，设计了一种简单而高效的解决方案。
该方案首先在无延迟环境中学习一个策略，然后在延迟环境中模仿该策略。
本文证明了在环境具有平滑性时，此类模仿策略相对于无延迟专家的性能界。
利用这些理论洞见，设计了 DIDA 算法，该算法遵循上述原则，采用 \textsc{DAgger} 作为模仿算法。
实验结果展示了 DIDA 相对于众多基线方法的有效性。
DIDA 几乎总能取得更优的回报，同时需要更少的样本。
有趣的是，本文展示了该算法如何经过修改适用于另外三种场景。
首先，当延迟为常值但非整数时，对 DIDA 进行简单调整即表现出良好的实验结果。
其次，虽然 DIDA 是在常值延迟上训练的，但本文为其在随机延迟上测试时提供了理论保证。
第三，将\Cref{chap:belief_based} 中的简单"技巧"应用于 DIDA，可利用针对某一延迟学习的策略高效地应用于更小延迟。
作为未来研究方向，可修改当前算法使其在随机延迟上训练，并学习\Cref{eq:belief_pol_stoch} 中的策略。
这样将获得比\Cref{thm:stoch_mdp_dida_bound} 提供更好的理论保证。
另一个方向是解决 DIDA 的主要缺点——对无延迟专家的依赖。
一种可能的思路是从延迟策略收集的数据集中离线学习一个无延迟专家。
这意味着不再需要从无延迟环境中采样任何样本。
